%==============================================================================
% DISCUSSION + LIMITATIONS (enxugada em 04/10/2026 para o camera-ready).
% GUARDRAIL: nenhuma limitacao foi REMOVIDA. As sete da versao avaliada
% continuam enunciadas (amostra pequena / sessao e robo unicos / contraste nao
% single-factor / estimadores closed-loop nao tentados / Ljung-Box / constantes
% cinematicas / rSim nao sintonizado + heading + condicao inicial) e entrou a
% que os revisores pediram (R1, R2): canais desacoplados, slip, saturacao e
% envelope dos logs. O rotulo "Limitations." em run-in custa zero linhas e
% mantem a subsecao visivel para o revisor. O que saiu: repeticoes de IV-D
% (os dois ajustes closed-loop, o limiar que engaja em 2 % das amostras) e o
% braco open-loop com translacao, que ja esta no Future Work.
% Para usar: em main.tex, trocar %==============================================================================
% DISCUSSION + LIMITATIONS — FUNDIDAS numa seção só (2026-08-10, pedido do
% autor; corte 9->6 páginas). A Seção VI separada foi eliminada e o arquivo
% 06_limitations.tex saiu do build.
% GUARDRAIL: nenhuma limitação foi REMOVIDA. Todas as sete continuam
% enunciadas, agora em um parágrafo denso: amostra pequena / sessão e robô
% únicos / contraste não single-factor / estimadores closed-loop não tentados /
% Ljung-Box / constantes cinemáticas / rSim não sintonizado + baseline ideal +
% mismatch de velocidade inicial / escopo de um único dataset e grade irregular.
% O detalhamento longo permanece em docs/BACKLOG_EXTENSAO_JINT.md item A.
%==============================================================================
\section{Discussion}
\label{sec:discussion}

\subsection{Identifiability is a property of the data, not the DOF}
\label{sec:discussion:inversion}
The central finding is not that yaw is harder to model than translation, but
that the \emph{same} structure and tooling give opposite verdicts depending on
the excitation regime: on closed-loop logs the yaw P1D meets none of the
acceptance criteria, while on a small dedicated collection it passes every
pre-registered gate with near-unity gain. Nothing in the pipeline distinguishes
the DOFs; only the data does, since those logs excite $v_x,v_y$ persistently
while the yaw command is a median $1\%$ of the per-wheel actuation.

The closed-loop gain estimate is not merely imprecise; it is wrong in the same
direction across two independent attempts, a direct high-rotation fit and a
regime-restricted linear fit. The dilution result of
Sec.~\ref{sec:results:omega} says what that estimator is actually fitting:
with the yaw command at about $1\%$ of the per-wheel setpoint, it explains only
$56\%$ of the variance of the measured rotation, the remainder plausibly coming
from slip, contact and traction asymmetry generated by the translational motion
(we measure the unexplained share, not its physical decomposition). Regressing
measured yaw on commanded yaw is then dominated by whatever correlates with the
command, and under feedback the controller issues yaw precisely in response to
those disturbances, which is the textbook closed-loop identification
hazard~\cite{gustavsson1977identification,vandenhof1998closedloop,forssell1999closedloop}.
This removes the reading that the actuation threshold simply switches yaw off,
since it engages in only ${\sim}2\%$ of samples, and replaces it with something
measurable. It does not isolate feedback from dilution, because the dedicated
collection removes both at once.

\subsection{What the model actually buys}
\label{sec:discussion:buys}
% REESCRITO 2026-08-10 (B4 da revisao V2): a leitura anterior ("quase tudo vem
% do lag") era artefato da ordem da ablacao cumulativa. O fatorial completo com
% Shapley mostra lag e dead time comparaveis.
A Simscape multibody twin of an SSL robot reports trajectory RMSE of
150--240~mm \cite{araujo2025digital}, against the $39$--$119$~mm reported here.
The two numbers are not comparable: that work scores whole trajectories from a
single start, so error accumulates without bound, whereas our windows are
re-anchored every $\SI{1.5}{s}$ and the metric is a median over windows. We
mention it only because the reader will otherwise make the comparison unaided.

The order-independent decomposition in Sec.~\ref{sec:results:rsim} says what the
model buys, and it is narrower than the fit statistics suggest. The identified
gain is almost irrelevant, which is expected once the gains sit near unity: what
matters is not \emph{how much} the robot moves per unit command but \emph{when}.
The time constant and the dead time split that credit in comparable parts, and
they are partly interchangeable, since both postpone the response. The practical
consequence is that a model of this family helps wherever commanded and attained
velocity differ by tens of milliseconds of lag, and helps little where they do
not. It also helps only in combination with a measured initial velocity: the
dynamics are what make that state usable, and the same decomposition shows
neither ingredient is sufficient alone.

\subsection{Limitations}
\label{sec:limitations}
% CORTE 2026-08-10 (9->6 pag.): parágrafo enxugado ~40%. GUARDRAIL mantido --
% todas as sete limitações continuam enunciadas, apenas sem detalhamento.
The yaw fit rests on six events from a single session, with no fine amplitude
sweep, so its threshold is an interval; all data come from one robot, one
surface and one hold-out split, and pre-registration and deterministic
re-execution mitigate but do not replace replication. The two yaw campaigns
differ in loop topology, yaw-to-translation ratio and session at once, so
separating the last two would need an open-loop arm with translation
superimposed, which we did not collect; nor did we attempt the closed-loop-aware
estimators (instrumental-variable, two-stage, joint input--output) the cited
literature proposes. Ljung--Box rejects residual whiteness for both yaw fits,
which the smoothing differentiator explains and which is why no parameter
intervals are reported. The firmware and CAD radius ratios still differ by
${\sim}3.3\%$. On evaluation, $99.6\%$ of hold-out windows are
translation-dominated, so rotational fidelity is not validated; the heading
advantage all but vanishes by $\SI{1.5}{s}$ (Table~\ref{tab:three}), so the
position claim does not transfer to heading in the same proportion; and the
rSim comparison remains fitted-versus-untuned and confounded by
% FIX 2026-08-10 (revisao V3, objecao 3a): esta frase ainda dizia que o e_ang do
% modelo EXCEDE o do rSim, contradizendo a Tabela III e o texto de IV-C, ja
% corrigidos quando o e_ang do rSim foi finalmente medido.
initial state, which the matched-delay ablation narrows but does not remove.
% source: results/identificacao_omega_20260724.md Etapa 4 ; qc_rotate_20260724_125828.md ;
%   config/robot_geometry.json ; docs/PREREGISTRO_OMEGA_20260724.md ;
%   results/window_errors_eang_summary_20260731_023856.md ;
%   results/eval_twin_vs_rsim_mode{A,B}_20260617_*.md

\subsection{A methodological lesson beyond this robot}
\label{sec:discussion:lesson}
% ELEVADO 2026-08-10 a procedimento nomeado (originalidade): antes era a moral
% de um resultado negativo; agora e uma triagem reutilizavel com regra de decisao
% e os tres DOFs deste robo como validacao.
Closed-loop operational data are the abundant resource in the SSL
\cite{melo2023dataset}, whereas a dedicated campaign requires robot access
outside routine play. The practical question is therefore not whether a log is
usable, but \emph{which channels} of it are --- and that can be answered before
fitting anything. We propose screening each degree of freedom on two quantities
computed from the log alone: the \emph{actuation share}, the fraction of the
commanded actuator effort attributable to that channel, and the \emph{explained
variance}, the fraction of the attained velocity its command accounts for once
aligned by the dead time. On this robot the two agree and separate the channels
cleanly: $v_x$ and $v_y$ sit at $94$--$96\%$ explained variance and dominate the
wheel setpoints, while $\omega$ sits at $56\%$ and carries $1\%$ of the
actuation. The rule we would apply is that a channel in the second regime should
prompt a short targeted collection rather than a more elaborate estimator, and
the screening costs nothing because both quantities come from data already
recorded.

This is not a restatement of persistent excitation, which is a spectral property
of the input \cite{ljung1999system}: a channel carrying $1\%$ of the actuation
could still be identifiable if its command were rich enough and the measurement
clean. The two are complementary. What the actuation share adds is that it
points at a physical cause --- here, a per-wheel threshold applied after the
command kinematics --- and therefore at a remedy, rather than only reporting
that an estimate came out badly.
 por
% %==============================================================================
% DISCUSSION + LIMITATIONS (enxugada em 04/10/2026 para o camera-ready).
% GUARDRAIL: nenhuma limitacao foi REMOVIDA. As sete da versao avaliada
% continuam enunciadas (amostra pequena / sessao e robo unicos / contraste nao
% single-factor / estimadores closed-loop nao tentados / Ljung-Box / constantes
% cinematicas / rSim nao sintonizado + heading + condicao inicial) e entrou a
% que os revisores pediram (R1, R2): canais desacoplados, slip, saturacao e
% envelope dos logs. O rotulo "Limitations." em run-in custa zero linhas e
% mantem a subsecao visivel para o revisor. O que saiu: repeticoes de IV-D
% (os dois ajustes closed-loop, o limiar que engaja em 2 % das amostras) e o
% braco open-loop com translacao, que ja esta no Future Work.
% Para usar: em main.tex, trocar %==============================================================================
% DISCUSSION + LIMITATIONS — FUNDIDAS numa seção só (2026-08-10, pedido do
% autor; corte 9->6 páginas). A Seção VI separada foi eliminada e o arquivo
% 06_limitations.tex saiu do build.
% GUARDRAIL: nenhuma limitação foi REMOVIDA. Todas as sete continuam
% enunciadas, agora em um parágrafo denso: amostra pequena / sessão e robô
% únicos / contraste não single-factor / estimadores closed-loop não tentados /
% Ljung-Box / constantes cinemáticas / rSim não sintonizado + baseline ideal +
% mismatch de velocidade inicial / escopo de um único dataset e grade irregular.
% O detalhamento longo permanece em docs/BACKLOG_EXTENSAO_JINT.md item A.
%==============================================================================
\section{Discussion}
\label{sec:discussion}

\subsection{Identifiability is a property of the data, not the DOF}
\label{sec:discussion:inversion}
The central finding is not that yaw is harder to model than translation, but
that the \emph{same} structure and tooling give opposite verdicts depending on
the excitation regime: on closed-loop logs the yaw P1D meets none of the
acceptance criteria, while on a small dedicated collection it passes every
pre-registered gate with near-unity gain. Nothing in the pipeline distinguishes
the DOFs; only the data does, since those logs excite $v_x,v_y$ persistently
while the yaw command is a median $1\%$ of the per-wheel actuation.

The closed-loop gain estimate is not merely imprecise; it is wrong in the same
direction across two independent attempts, a direct high-rotation fit and a
regime-restricted linear fit. The dilution result of
Sec.~\ref{sec:results:omega} says what that estimator is actually fitting:
with the yaw command at about $1\%$ of the per-wheel setpoint, it explains only
$56\%$ of the variance of the measured rotation, the remainder plausibly coming
from slip, contact and traction asymmetry generated by the translational motion
(we measure the unexplained share, not its physical decomposition). Regressing
measured yaw on commanded yaw is then dominated by whatever correlates with the
command, and under feedback the controller issues yaw precisely in response to
those disturbances, which is the textbook closed-loop identification
hazard~\cite{gustavsson1977identification,vandenhof1998closedloop,forssell1999closedloop}.
This removes the reading that the actuation threshold simply switches yaw off,
since it engages in only ${\sim}2\%$ of samples, and replaces it with something
measurable. It does not isolate feedback from dilution, because the dedicated
collection removes both at once.

\subsection{What the model actually buys}
\label{sec:discussion:buys}
% REESCRITO 2026-08-10 (B4 da revisao V2): a leitura anterior ("quase tudo vem
% do lag") era artefato da ordem da ablacao cumulativa. O fatorial completo com
% Shapley mostra lag e dead time comparaveis.
A Simscape multibody twin of an SSL robot reports trajectory RMSE of
150--240~mm \cite{araujo2025digital}, against the $39$--$119$~mm reported here.
The two numbers are not comparable: that work scores whole trajectories from a
single start, so error accumulates without bound, whereas our windows are
re-anchored every $\SI{1.5}{s}$ and the metric is a median over windows. We
mention it only because the reader will otherwise make the comparison unaided.

The order-independent decomposition in Sec.~\ref{sec:results:rsim} says what the
model buys, and it is narrower than the fit statistics suggest. The identified
gain is almost irrelevant, which is expected once the gains sit near unity: what
matters is not \emph{how much} the robot moves per unit command but \emph{when}.
The time constant and the dead time split that credit in comparable parts, and
they are partly interchangeable, since both postpone the response. The practical
consequence is that a model of this family helps wherever commanded and attained
velocity differ by tens of milliseconds of lag, and helps little where they do
not. It also helps only in combination with a measured initial velocity: the
dynamics are what make that state usable, and the same decomposition shows
neither ingredient is sufficient alone.

\subsection{Limitations}
\label{sec:limitations}
% CORTE 2026-08-10 (9->6 pag.): parágrafo enxugado ~40%. GUARDRAIL mantido --
% todas as sete limitações continuam enunciadas, apenas sem detalhamento.
The yaw fit rests on six events from a single session, with no fine amplitude
sweep, so its threshold is an interval; all data come from one robot, one
surface and one hold-out split, and pre-registration and deterministic
re-execution mitigate but do not replace replication. The two yaw campaigns
differ in loop topology, yaw-to-translation ratio and session at once, so
separating the last two would need an open-loop arm with translation
superimposed, which we did not collect; nor did we attempt the closed-loop-aware
estimators (instrumental-variable, two-stage, joint input--output) the cited
literature proposes. Ljung--Box rejects residual whiteness for both yaw fits,
which the smoothing differentiator explains and which is why no parameter
intervals are reported. The firmware and CAD radius ratios still differ by
${\sim}3.3\%$. On evaluation, $99.6\%$ of hold-out windows are
translation-dominated, so rotational fidelity is not validated; the heading
advantage all but vanishes by $\SI{1.5}{s}$ (Table~\ref{tab:three}), so the
position claim does not transfer to heading in the same proportion; and the
rSim comparison remains fitted-versus-untuned and confounded by
% FIX 2026-08-10 (revisao V3, objecao 3a): esta frase ainda dizia que o e_ang do
% modelo EXCEDE o do rSim, contradizendo a Tabela III e o texto de IV-C, ja
% corrigidos quando o e_ang do rSim foi finalmente medido.
initial state, which the matched-delay ablation narrows but does not remove.
% source: results/identificacao_omega_20260724.md Etapa 4 ; qc_rotate_20260724_125828.md ;
%   config/robot_geometry.json ; docs/PREREGISTRO_OMEGA_20260724.md ;
%   results/window_errors_eang_summary_20260731_023856.md ;
%   results/eval_twin_vs_rsim_mode{A,B}_20260617_*.md

\subsection{A methodological lesson beyond this robot}
\label{sec:discussion:lesson}
% ELEVADO 2026-08-10 a procedimento nomeado (originalidade): antes era a moral
% de um resultado negativo; agora e uma triagem reutilizavel com regra de decisao
% e os tres DOFs deste robo como validacao.
Closed-loop operational data are the abundant resource in the SSL
\cite{melo2023dataset}, whereas a dedicated campaign requires robot access
outside routine play. The practical question is therefore not whether a log is
usable, but \emph{which channels} of it are --- and that can be answered before
fitting anything. We propose screening each degree of freedom on two quantities
computed from the log alone: the \emph{actuation share}, the fraction of the
commanded actuator effort attributable to that channel, and the \emph{explained
variance}, the fraction of the attained velocity its command accounts for once
aligned by the dead time. On this robot the two agree and separate the channels
cleanly: $v_x$ and $v_y$ sit at $94$--$96\%$ explained variance and dominate the
wheel setpoints, while $\omega$ sits at $56\%$ and carries $1\%$ of the
actuation. The rule we would apply is that a channel in the second regime should
prompt a short targeted collection rather than a more elaborate estimator, and
the screening costs nothing because both quantities come from data already
recorded.

This is not a restatement of persistent excitation, which is a spectral property
of the input \cite{ljung1999system}: a channel carrying $1\%$ of the actuation
could still be identifiable if its command were rich enough and the measurement
clean. The two are complementary. What the actuation share adds is that it
points at a physical cause --- here, a per-wheel threshold applied after the
command kinematics --- and therefore at a remedy, rather than only reporting
that an estimate came out badly.
 por
% %==============================================================================
% DISCUSSION + LIMITATIONS (enxugada em 04/10/2026 para o camera-ready).
% GUARDRAIL: nenhuma limitacao foi REMOVIDA. As sete da versao avaliada
% continuam enunciadas (amostra pequena / sessao e robo unicos / contraste nao
% single-factor / estimadores closed-loop nao tentados / Ljung-Box / constantes
% cinematicas / rSim nao sintonizado + heading + condicao inicial) e entrou a
% que os revisores pediram (R1, R2): canais desacoplados, slip, saturacao e
% envelope dos logs. O rotulo "Limitations." em run-in custa zero linhas e
% mantem a subsecao visivel para o revisor. O que saiu: repeticoes de IV-D
% (os dois ajustes closed-loop, o limiar que engaja em 2 % das amostras) e o
% braco open-loop com translacao, que ja esta no Future Work.
% Para usar: em main.tex, trocar %==============================================================================
% DISCUSSION + LIMITATIONS — FUNDIDAS numa seção só (2026-08-10, pedido do
% autor; corte 9->6 páginas). A Seção VI separada foi eliminada e o arquivo
% 06_limitations.tex saiu do build.
% GUARDRAIL: nenhuma limitação foi REMOVIDA. Todas as sete continuam
% enunciadas, agora em um parágrafo denso: amostra pequena / sessão e robô
% únicos / contraste não single-factor / estimadores closed-loop não tentados /
% Ljung-Box / constantes cinemáticas / rSim não sintonizado + baseline ideal +
% mismatch de velocidade inicial / escopo de um único dataset e grade irregular.
% O detalhamento longo permanece em docs/BACKLOG_EXTENSAO_JINT.md item A.
%==============================================================================
\section{Discussion}
\label{sec:discussion}

\subsection{Identifiability is a property of the data, not the DOF}
\label{sec:discussion:inversion}
The central finding is not that yaw is harder to model than translation, but
that the \emph{same} structure and tooling give opposite verdicts depending on
the excitation regime: on closed-loop logs the yaw P1D meets none of the
acceptance criteria, while on a small dedicated collection it passes every
pre-registered gate with near-unity gain. Nothing in the pipeline distinguishes
the DOFs; only the data does, since those logs excite $v_x,v_y$ persistently
while the yaw command is a median $1\%$ of the per-wheel actuation.

The closed-loop gain estimate is not merely imprecise; it is wrong in the same
direction across two independent attempts, a direct high-rotation fit and a
regime-restricted linear fit. The dilution result of
Sec.~\ref{sec:results:omega} says what that estimator is actually fitting:
with the yaw command at about $1\%$ of the per-wheel setpoint, it explains only
$56\%$ of the variance of the measured rotation, the remainder plausibly coming
from slip, contact and traction asymmetry generated by the translational motion
(we measure the unexplained share, not its physical decomposition). Regressing
measured yaw on commanded yaw is then dominated by whatever correlates with the
command, and under feedback the controller issues yaw precisely in response to
those disturbances, which is the textbook closed-loop identification
hazard~\cite{gustavsson1977identification,vandenhof1998closedloop,forssell1999closedloop}.
This removes the reading that the actuation threshold simply switches yaw off,
since it engages in only ${\sim}2\%$ of samples, and replaces it with something
measurable. It does not isolate feedback from dilution, because the dedicated
collection removes both at once.

\subsection{What the model actually buys}
\label{sec:discussion:buys}
% REESCRITO 2026-08-10 (B4 da revisao V2): a leitura anterior ("quase tudo vem
% do lag") era artefato da ordem da ablacao cumulativa. O fatorial completo com
% Shapley mostra lag e dead time comparaveis.
A Simscape multibody twin of an SSL robot reports trajectory RMSE of
150--240~mm \cite{araujo2025digital}, against the $39$--$119$~mm reported here.
The two numbers are not comparable: that work scores whole trajectories from a
single start, so error accumulates without bound, whereas our windows are
re-anchored every $\SI{1.5}{s}$ and the metric is a median over windows. We
mention it only because the reader will otherwise make the comparison unaided.

The order-independent decomposition in Sec.~\ref{sec:results:rsim} says what the
model buys, and it is narrower than the fit statistics suggest. The identified
gain is almost irrelevant, which is expected once the gains sit near unity: what
matters is not \emph{how much} the robot moves per unit command but \emph{when}.
The time constant and the dead time split that credit in comparable parts, and
they are partly interchangeable, since both postpone the response. The practical
consequence is that a model of this family helps wherever commanded and attained
velocity differ by tens of milliseconds of lag, and helps little where they do
not. It also helps only in combination with a measured initial velocity: the
dynamics are what make that state usable, and the same decomposition shows
neither ingredient is sufficient alone.

\subsection{Limitations}
\label{sec:limitations}
% CORTE 2026-08-10 (9->6 pag.): parágrafo enxugado ~40%. GUARDRAIL mantido --
% todas as sete limitações continuam enunciadas, apenas sem detalhamento.
The yaw fit rests on six events from a single session, with no fine amplitude
sweep, so its threshold is an interval; all data come from one robot, one
surface and one hold-out split, and pre-registration and deterministic
re-execution mitigate but do not replace replication. The two yaw campaigns
differ in loop topology, yaw-to-translation ratio and session at once, so
separating the last two would need an open-loop arm with translation
superimposed, which we did not collect; nor did we attempt the closed-loop-aware
estimators (instrumental-variable, two-stage, joint input--output) the cited
literature proposes. Ljung--Box rejects residual whiteness for both yaw fits,
which the smoothing differentiator explains and which is why no parameter
intervals are reported. The firmware and CAD radius ratios still differ by
${\sim}3.3\%$. On evaluation, $99.6\%$ of hold-out windows are
translation-dominated, so rotational fidelity is not validated; the heading
advantage all but vanishes by $\SI{1.5}{s}$ (Table~\ref{tab:three}), so the
position claim does not transfer to heading in the same proportion; and the
rSim comparison remains fitted-versus-untuned and confounded by
% FIX 2026-08-10 (revisao V3, objecao 3a): esta frase ainda dizia que o e_ang do
% modelo EXCEDE o do rSim, contradizendo a Tabela III e o texto de IV-C, ja
% corrigidos quando o e_ang do rSim foi finalmente medido.
initial state, which the matched-delay ablation narrows but does not remove.
% source: results/identificacao_omega_20260724.md Etapa 4 ; qc_rotate_20260724_125828.md ;
%   config/robot_geometry.json ; docs/PREREGISTRO_OMEGA_20260724.md ;
%   results/window_errors_eang_summary_20260731_023856.md ;
%   results/eval_twin_vs_rsim_mode{A,B}_20260617_*.md

\subsection{A methodological lesson beyond this robot}
\label{sec:discussion:lesson}
% ELEVADO 2026-08-10 a procedimento nomeado (originalidade): antes era a moral
% de um resultado negativo; agora e uma triagem reutilizavel com regra de decisao
% e os tres DOFs deste robo como validacao.
Closed-loop operational data are the abundant resource in the SSL
\cite{melo2023dataset}, whereas a dedicated campaign requires robot access
outside routine play. The practical question is therefore not whether a log is
usable, but \emph{which channels} of it are --- and that can be answered before
fitting anything. We propose screening each degree of freedom on two quantities
computed from the log alone: the \emph{actuation share}, the fraction of the
commanded actuator effort attributable to that channel, and the \emph{explained
variance}, the fraction of the attained velocity its command accounts for once
aligned by the dead time. On this robot the two agree and separate the channels
cleanly: $v_x$ and $v_y$ sit at $94$--$96\%$ explained variance and dominate the
wheel setpoints, while $\omega$ sits at $56\%$ and carries $1\%$ of the
actuation. The rule we would apply is that a channel in the second regime should
prompt a short targeted collection rather than a more elaborate estimator, and
the screening costs nothing because both quantities come from data already
recorded.

This is not a restatement of persistent excitation, which is a spectral property
of the input \cite{ljung1999system}: a channel carrying $1\%$ of the actuation
could still be identifiable if its command were rich enough and the measurement
clean. The two are complementary. What the actuation share adds is that it
points at a physical cause --- here, a per-wheel threshold applied after the
command kinematics --- and therefore at a remedy, rather than only reporting
that an estimate came out badly.
 por
% %==============================================================================
% DISCUSSION + LIMITATIONS (enxugada em 04/10/2026 para o camera-ready).
% GUARDRAIL: nenhuma limitacao foi REMOVIDA. As sete da versao avaliada
% continuam enunciadas (amostra pequena / sessao e robo unicos / contraste nao
% single-factor / estimadores closed-loop nao tentados / Ljung-Box / constantes
% cinematicas / rSim nao sintonizado + heading + condicao inicial) e entrou a
% que os revisores pediram (R1, R2): canais desacoplados, slip, saturacao e
% envelope dos logs. O rotulo "Limitations." em run-in custa zero linhas e
% mantem a subsecao visivel para o revisor. O que saiu: repeticoes de IV-D
% (os dois ajustes closed-loop, o limiar que engaja em 2 % das amostras) e o
% braco open-loop com translacao, que ja esta no Future Work.
% Para usar: em main.tex, trocar \input{sections/05_discussion} por
% \input{sections/05_discussion_enxuta}, ou copiar este conteudo por cima.
% Fontes: results/identificacao_omega_20260724.md Etapa 4 ;
%   results/qc_rotate_20260724_125828.md ; config/robot_geometry.json ;
%   docs/PREREGISTRO_OMEGA_20260724.md ; results/window_errors_eang_summary_20260731_023856.md ;
%   results/eval_twin_vs_rsim_mode{A,B}_20260617_*.md ; results/eval_rsim_delay_ablation_20260809_2154.md ;
%   config/robot_limits.json (planta publicada sem saturacao nem a_max).
%==============================================================================
\section{Discussion}
\label{sec:discussion}

The central finding is that the \emph{same} structure and tooling give opposite verdicts depending on the excitation regime: on closed-loop logs the yaw P1D meets none of the acceptance criteria, while on a small dedicated collection it passes every pre-registered gate with near-unity gain. Only the data distinguishes the DOFs, and the dilution result of Sec.~\ref{sec:results:omega} says what the closed-loop estimator is actually fitting: with the yaw command at about $1\%$ of the per-wheel setpoint, it explains only $56\%$ of the measured rotation, the remainder plausibly coming from slip, contact and traction asymmetry generated by the translational motion. Under feedback the controller issues yaw precisely in response to those disturbances, so the regression is dominated by whatever correlates with the command, the textbook closed-loop identification hazard~\cite{gustavsson1977identification,vandenhof1998closedloop,forssell1999closedloop}; the dedicated collection removes dilution and feedback at once, so the two are not isolated here.

\textit{Limitations.} The model is deliberately simple: its three channels are decoupled, so the translation-induced yaw above is exactly what it cannot represent and what a coupled (MIMO) identification would have to capture; it is linear, with no slip state, saturation or acceleration limit, and it is fitted and scored inside the envelope of these logs, so its validity in aggressive or mixed-motion regimes is unverified. The yaw fit rests on six events from one session without a fine amplitude sweep, so its threshold is an interval; all data come from one robot, one surface and one hold-out split, and pre-registration and deterministic re-execution mitigate but do not replace replication. The two yaw campaigns differ in loop topology, yaw-to-translation ratio and session at once, and closed-loop-aware estimators were not attempted. Ljung--Box rejects residual whiteness for both yaw fits, as the smoothing differentiator predicts, so no parameter intervals are reported; the firmware and CAD radius ratios differ by ${\sim}3.3\%$. Finally, $99.6\%$ of hold-out windows are translation-dominated, so rotational fidelity is not validated, the heading advantage all but vanishes by $\SI{1.5}{s}$ (Table~\ref{tab:three}), and the rSim comparison remains fitted-versus-untuned and confounded by initial state, which the matched-delay ablation narrows but does not remove.
, ou copiar este conteudo por cima.
% Fontes: results/identificacao_omega_20260724.md Etapa 4 ;
%   results/qc_rotate_20260724_125828.md ; config/robot_geometry.json ;
%   docs/PREREGISTRO_OMEGA_20260724.md ; results/window_errors_eang_summary_20260731_023856.md ;
%   results/eval_twin_vs_rsim_mode{A,B}_20260617_*.md ; results/eval_rsim_delay_ablation_20260809_2154.md ;
%   config/robot_limits.json (planta publicada sem saturacao nem a_max).
%==============================================================================
\section{Discussion}
\label{sec:discussion}

The central finding is that the \emph{same} structure and tooling give opposite verdicts depending on the excitation regime: on closed-loop logs the yaw P1D meets none of the acceptance criteria, while on a small dedicated collection it passes every pre-registered gate with near-unity gain. Only the data distinguishes the DOFs, and the dilution result of Sec.~\ref{sec:results:omega} says what the closed-loop estimator is actually fitting: with the yaw command at about $1\%$ of the per-wheel setpoint, it explains only $56\%$ of the measured rotation, the remainder plausibly coming from slip, contact and traction asymmetry generated by the translational motion. Under feedback the controller issues yaw precisely in response to those disturbances, so the regression is dominated by whatever correlates with the command, the textbook closed-loop identification hazard~\cite{gustavsson1977identification,vandenhof1998closedloop,forssell1999closedloop}; the dedicated collection removes dilution and feedback at once, so the two are not isolated here.

\textit{Limitations.} The model is deliberately simple: its three channels are decoupled, so the translation-induced yaw above is exactly what it cannot represent and what a coupled (MIMO) identification would have to capture; it is linear, with no slip state, saturation or acceleration limit, and it is fitted and scored inside the envelope of these logs, so its validity in aggressive or mixed-motion regimes is unverified. The yaw fit rests on six events from one session without a fine amplitude sweep, so its threshold is an interval; all data come from one robot, one surface and one hold-out split, and pre-registration and deterministic re-execution mitigate but do not replace replication. The two yaw campaigns differ in loop topology, yaw-to-translation ratio and session at once, and closed-loop-aware estimators were not attempted. Ljung--Box rejects residual whiteness for both yaw fits, as the smoothing differentiator predicts, so no parameter intervals are reported; the firmware and CAD radius ratios differ by ${\sim}3.3\%$. Finally, $99.6\%$ of hold-out windows are translation-dominated, so rotational fidelity is not validated, the heading advantage all but vanishes by $\SI{1.5}{s}$ (Table~\ref{tab:three}), and the rSim comparison remains fitted-versus-untuned and confounded by initial state, which the matched-delay ablation narrows but does not remove.
, ou copiar este conteudo por cima.
% Fontes: results/identificacao_omega_20260724.md Etapa 4 ;
%   results/qc_rotate_20260724_125828.md ; config/robot_geometry.json ;
%   docs/PREREGISTRO_OMEGA_20260724.md ; results/window_errors_eang_summary_20260731_023856.md ;
%   results/eval_twin_vs_rsim_mode{A,B}_20260617_*.md ; results/eval_rsim_delay_ablation_20260809_2154.md ;
%   config/robot_limits.json (planta publicada sem saturacao nem a_max).
%==============================================================================
\section{Discussion}
\label{sec:discussion}

The central finding is that the \emph{same} structure and tooling give opposite verdicts depending on the excitation regime: on closed-loop logs the yaw P1D meets none of the acceptance criteria, while on a small dedicated collection it passes every pre-registered gate with near-unity gain. Only the data distinguishes the DOFs, and the dilution result of Sec.~\ref{sec:results:omega} says what the closed-loop estimator is actually fitting: with the yaw command at about $1\%$ of the per-wheel setpoint, it explains only $56\%$ of the measured rotation, the remainder plausibly coming from slip, contact and traction asymmetry generated by the translational motion. Under feedback the controller issues yaw precisely in response to those disturbances, so the regression is dominated by whatever correlates with the command, the textbook closed-loop identification hazard~\cite{gustavsson1977identification,vandenhof1998closedloop,forssell1999closedloop}; the dedicated collection removes dilution and feedback at once, so the two are not isolated here.

\textit{Limitations.} The model is deliberately simple: its three channels are decoupled, so the translation-induced yaw above is exactly what it cannot represent and what a coupled (MIMO) identification would have to capture; it is linear, with no slip state, saturation or acceleration limit, and it is fitted and scored inside the envelope of these logs, so its validity in aggressive or mixed-motion regimes is unverified. The yaw fit rests on six events from one session without a fine amplitude sweep, so its threshold is an interval; all data come from one robot, one surface and one hold-out split, and pre-registration and deterministic re-execution mitigate but do not replace replication. The two yaw campaigns differ in loop topology, yaw-to-translation ratio and session at once, and closed-loop-aware estimators were not attempted. Ljung--Box rejects residual whiteness for both yaw fits, as the smoothing differentiator predicts, so no parameter intervals are reported; the firmware and CAD radius ratios differ by ${\sim}3.3\%$. Finally, $99.6\%$ of hold-out windows are translation-dominated, so rotational fidelity is not validated, the heading advantage all but vanishes by $\SI{1.5}{s}$ (Table~\ref{tab:three}), and the rSim comparison remains fitted-versus-untuned and confounded by initial state, which the matched-delay ablation narrows but does not remove.
, ou copiar este conteudo por cima.
% Fontes: results/identificacao_omega_20260724.md Etapa 4 ;
%   results/qc_rotate_20260724_125828.md ; config/robot_geometry.json ;
%   docs/PREREGISTRO_OMEGA_20260724.md ; results/window_errors_eang_summary_20260731_023856.md ;
%   results/eval_twin_vs_rsim_mode{A,B}_20260617_*.md ; results/eval_rsim_delay_ablation_20260809_2154.md ;
%   config/robot_limits.json (planta publicada sem saturacao nem a_max).
%==============================================================================
\section{Discussion}
\label{sec:discussion}

The central finding is that the \emph{same} structure and tooling give opposite verdicts depending on the excitation regime: on closed-loop logs the yaw P1D meets none of the acceptance criteria, while on a small dedicated collection it passes every pre-registered gate with near-unity gain. Only the data distinguishes the DOFs, and the dilution result of Sec.~\ref{sec:results:omega} says what the closed-loop estimator is actually fitting: with the yaw command at about $1\%$ of the per-wheel setpoint, it explains only $56\%$ of the measured rotation, the remainder plausibly coming from slip, contact and traction asymmetry generated by the translational motion. Under feedback the controller issues yaw precisely in response to those disturbances, so the regression is dominated by whatever correlates with the command, the textbook closed-loop identification hazard~\cite{gustavsson1977identification,vandenhof1998closedloop,forssell1999closedloop}; the dedicated collection removes dilution and feedback at once, so the two are not isolated here.

\textit{Limitations.} The model is deliberately simple: its three channels are decoupled, so the translation-induced yaw above is exactly what it cannot represent and what a coupled (MIMO) identification would have to capture; it is linear, with no slip state, saturation or acceleration limit, and it is fitted and scored inside the envelope of these logs, so its validity in aggressive or mixed-motion regimes is unverified. The yaw fit rests on six events from one session without a fine amplitude sweep, so its threshold is an interval; all data come from one robot, one surface and one hold-out split, and pre-registration and deterministic re-execution mitigate but do not replace replication. The two yaw campaigns differ in loop topology, yaw-to-translation ratio and session at once, and closed-loop-aware estimators were not attempted. Ljung--Box rejects residual whiteness for both yaw fits, as the smoothing differentiator predicts, so no parameter intervals are reported; the firmware and CAD radius ratios differ by ${\sim}3.3\%$. Finally, $99.6\%$ of hold-out windows are translation-dominated, so rotational fidelity is not validated, the heading advantage all but vanishes by $\SI{1.5}{s}$ (Table~\ref{tab:three}), and the rSim comparison remains fitted-versus-untuned and confounded by initial state, which the matched-delay ablation narrows but does not remove.
